\section{Unchanged labels with only a finite second moment}
\label{sec:general-shallow-finitevariance}
The sharp Gaussian shallow oracle also holds under the sole label
assumption $\mathbb E Y^2<\infty$. This is a high-probability result,
not an expectation bound. No labels or network outputs are clipped.

\begin{lemma}[A square-integrable multiplier]
Let $a(X,Y)$ be any fixed square-integrable measurable function.
For the normalized ReLU atoms of the preceding section,
\[
 \mathbb E_S\left[
 \sup_{u\in\mathcal H_d}
 |(\mathbb P_n-\mathbb P)[a(X,Y)u(X)]|^2\right]
 \le \frac{C\mathbb E a(X,Y)^2}{n}.                      \tag{GF1}
\]
\end{lemma}
\begin{proof}
An independent ghost sample and Jensen's inequality symmetrize
the squared supremum, with a numerical factor, to
\[
 \mathbb E_{S,\epsilon}
 \left[\sup_{u\in\mathcal H_d}
       \left|n^{-1}\sum_i\epsilon_i a_i u(X_i)\right|^2\right].
\]
Conditional on the sample, the mean of this signed supremum is at
most $C\sqrt{\sum_i a_i^2}/n$, by the weighted contraction argument
above. Changing one Rademacher sign changes the supremum by at most
$2\sqrt2|a_i|/n$. The bounded-differences variance inequality therefore
bounds its conditional variance by $C\sum_i a_i^2/n^2$.
Adding the squared conditional mean, then averaging, proves (GF1).
Truncation of $a$ followed by the same $L^2$ estimate justifies every
step for general square-integrable multipliers.
\end{proof}

\begin{theorem}[A finite-variance, teacher-independent likelihood schedule]
\label{thm:general-shallow-finitevariance}
Let inputs have any law on $[-1,1]^d$, and suppose
$\mathbb E Y^2\le M_2<\infty$. Fix $0<\delta<1/2$ and set
\begin{align*}
 L&=\frac{8(M_2+1)}\delta,\\
 b&=\frac{c_0\log^4(\mathrm e(n+1)(d+L+1)/\delta)}n,\\
 \lambda&=c_1\{\sqrt{Lb}+(d+2)b\}.
\end{align*}
Assume $\lambda\le1$, with sufficiently large fixed numerical
$c_0,c_1$. For any deterministic finite shallow ReLU comparison $g$,
with probability at least $1-\delta$,
\[
 R(\bar z_\lambda)\le\Pi_\lambda R(z)
 \le R(g)+C\lambda
       [P_g+1+(d+2)\log((n+1)(d+L+1))].                 \tag{GF2}
\]
This event is common to every $N\ge n$ and $0<\nu\le1$.
The prior is the original independent Gaussian prior; the likelihood
uses the actual raw observations, and its physical variance is
$n\lambda\nu/N$. The deterministic schedule depends on the declared
second-moment bound and confidence, but not on the comparator.
\end{theorem}
\begin{proof}
Markov's inequality gives $\mathbb P_nY^2\le L$ except with
probability at most $\delta/8$. On this event the moment version
(GM5) of the empirical midpoint inequality applies, and
\[
 Z\ge\exp[-\omega(L+1)/(2\lambda)],\qquad
 \Pi Q^2\le C[(d+2)\nu+(L+1)/\lambda]^2.
\]
The proof of (GS4), with this evidence bound, gives
\[
 \Pi[(1+Q^2)\mathbf1_{Q>q_*}]\le n^{-40},\qquad
 q_*=A\{d+2+(L+1)/\lambda\}.                           \tag{GF3}
\]
This is a deterministic assertion on the empirical-second-moment
event, simultaneous in width and prior scale.

First consider $R(g)\le R(0)+1\le M_2+1$ and
$\lambda(P_g+1)\le cL$. Apply (GF1) to $a=g-Y$ and Markov's
inequality. Except with probability $\delta/4$, its multiplier
supremum is at most $C\sqrt{(M_2+1)/(n\delta)}\le C\sqrt{Lb}$.
Intersect this with the square-norm event (GS3), with its confidence
allocation reduced by a fixed factor. Expanding the loss difference
gives, uniformly on $Q\le q_*$,
\begin{align*}
 R(z)-R(g)\le{}&\widehat R(z)-\widehat R(g)
       +\tfrac14\mathbb P_n(z-g)^2\\
       &+Cb(P+P_g+1)^2+C(P+P_g)\sqrt{Lb}.
\end{align*}
Now $bq_*\le c'\lambda$ follows from
$\lambda\ge c_1\sqrt{Lb}$ and $\lambda\ge c_1(d+2)b$.
Also $b(P_g+1)^2\le C\lambda(P_g+1)$ in the stated comparison
range, since $\lambda^2\ge c_1^2Lb$.
Enlarging $c_1$ consequently bounds the last line by
$\lambda P/2+C\lambda(P_g+1)$.

Integrate, apply (GM5), and use (GF3) for the discarded tail.
The empirical comparator risk is at most $C(P_g+1)^2+2L$.
Since $\lambda\le1$ forces $L,d\le Cn$ and the nontrivial branch
forces $P_g\le Cn$, these tail corrections are negligible.
This proves (GF2) in that branch.

Apply the same argument to zero to get
\[
 \Pi R(z)\le R(0)+C\lambda
              [1+(d+2)\log((n+1)(d+L+1))].
\]
If $R(g)>R(0)$ this already suffices. If
$\lambda(P_g+1)>cL$, the excess resource term can dominate
$R(0)\le M_2\le L$, after enlarging its numerical constant.
Thus all comparisons are covered. Jensen proves the mean bound.
\end{proof}

\paragraph{No second-moment bound is needed asymptotically.}
Fix input dimension and a confidence $0<\delta<1/2$. The one
deterministic schedule
\[
 \lambda_n=\frac{\log^6(n+2)}{\sqrt n}
                  +(d+2)\frac{\log^6(n+2)}n              \tag{GF4}
\]
eventually satisfies every lower bound in the proof for each fixed
finite $M_2$. Its onset may depend on that unknown label law and on
$\delta$; the quantitative finite-sample theorem above uses a
declared moment bound. At fixed coefficient cap $B$ and dimension,
the raw posterior at (GF4) therefore competes with width
$C_n=\lfloor n^{1/4}\rfloor$ shallow ReLU comparisons with
high-probability gap $\widetilde O(n^{-1/4})$, at every $N\ge n$
and every mean-field-or-stronger scale. More generally any prescribed
$C_nB_n^2=o(\sqrt n/\log^7 n)$ gives a vanishing bound. This is a
bounded-coefficient comparison, not a coefficient-free theorem on
arbitrary input laws.

\paragraph{Fixed classes and optional simultaneous budgets.}
For each declared finite coefficient/width class choose a
deterministic population minimizer, which exists by continuity on
its compact parameter set. This gives its oracle immediately.
The fixed-class result is sufficient for the growing class in (GF4).
One can also keep one event over polynomially many nontrivial budget
classes: apply (GF1) only to the common multiplier $a=Y$, and handle
each bounded $g$ multiplier separately by Bousquet's inequality.
Since $\mathbb E g^2\le C(M_2+1)$ for the chosen near-minimizers,
the resulting simultaneous multiplier coefficient is
$C\{\sqrt{(M_2+1)b}+(P_g+1)b+\sqrt{(M_2+1)/(n\delta)}\}$.
It obeys the same absorption inequalities. A polynomial finite union
changes only the logarithms already included in $b$; larger budgets
are again handled by zero. No union of Markov events over all classes
is needed.

\begin{corollary}[Universal risk consistency in probability]
\label{cor:gaussian-universal-consistency}
Fix the input dimension and any probability law on the input cube,
and assume only $\mathbb E Y^2<\infty$. For any prescribed Gaussian scales $0<\nu_n\le1$ and finite widths
$N_n\to\infty$, however slowly, the ordinary raw posterior mean is
risk-consistent in probability under either of two choices: a shallow
ReLU learner at the universal $\log^8$ schedule of
Section~\ref{sec:all-width-relu}, or a bounded, nonconstant,
nondecreasing Lipschitz activation at $\lambda_n=n^{-3/4}$.
The posterior-average risk has both conclusions. Neither schedule
needs the unknown label second moment.
\end{corollary}
\begin{proof}
Finite shallow ridge networks with any of these continuous
nonpolynomial activations are uniformly dense in continuous functions
on the cube~\cite[Theorem~1]{leshno}. Continuous functions are dense in
$L^2(P_X)$, and the regression function is square-integrable.
Consequently, for every $\varepsilon>0$, there is a fixed finite
comparison network $g$ with
$R(g)\le R_*+\varepsilon$, where
$R_*=\mathbb E\operatorname{Var}(Y\mid X)$.
For this fixed $g$, the ReLU remainder in (GA3) at the universal
schedule of Section~\ref{sec:all-width-relu} tends to zero;
in the bounded-activation case $\mathcal A_n(g)=O(\log n)$ and (AW3)
at $\lambda_n=n^{-3/4}$ also has a vanishing remainder. In both cases the
width requirement is eventually met because the fixed comparison
has finitely many neurons and $N_n\to\infty$.
Both statements hold for every fixed confidence level and uniformly
over the eligible prescribed widths and scales. Since test risk is
at least $R_*$, first choose the approximation error small and then
let the confidence failure probability tend to zero. This proves
convergence in probability. It does not assert expected-risk
consistency for arbitrary finite-variance labels. This is a standard
approximation consequence of the quantitative oracles.
\end{proof}
